\begin{resumo}[Abstract]
 \begin{otherlanguage*}{english}
This work presents the proposal and initial development of a Public Services
Evaluator for digital public services, in the context of the digital
transformation of the Brazilian State and the need to make citizen experience
evaluation more useful for continuous improvement. The research problem is
related to the low usefulness of evaluations displayed only at the end of the
user journey, since many citizens face difficulties before completing a service
and do not always answer traditional forms. The general objective is to propose
and prototype a solution capable of recording navigation difficulty signals,
suggesting redirection in priority services, requesting low-effort contextual
feedback, and presenting aggregated indicators to support digital public service
improvement. The methodology combined bibliographic and normative review, Design
Science Research, incremental prototyping, and Lean Inception. Based on product
vision, scope delimitation, product goals, personas, journeys, feature
brainstorming, technical, business and UX review, sequencing, and MVP Canvas
activities, the MVP was defined with anonymous technical sessions, click and time
tracking, simple rule-based detection, contextual feedback, a comparative static
form, and an administrative dashboard. The prototype was implemented with a local
API, simulated public portal pages, a feedback form with stars, attributes, NPS
and comments, and SQLite data storage. As a result, the work demonstrates the
conceptual and technical feasibility of a solution that transforms navigation
events and declared feedback into administrative evidence. It concludes that
contextual feedback can complement official evaluation instruments, provided that
it is used for service improvement, respects privacy, includes institutional data
governance, and is later validated with real users.

   \vspace{\onelineskip}

   \noindent
   \textbf{Key-words}: public services. contextual feedback. digital government.
   user experience. prototype.
 \end{otherlanguage*}
\end{resumo}
